\PassOptionsToPackage{dvipsnames,svgnames,table}{xcolor}
\documentclass[10pt]{article}
\usepackage[a4paper,margin=16mm]{geometry}
\usepackage[no-math]{fontspec}\setmainfont{Latin Modern Roman}
\usepackage{amsmath,amssymb,amsthm,mathtools,graphicx,xcolor,hyperref,xurl}
\usepackage[shortlabels]{enumitem}
\setlength{\emergencystretch}{4em}
\newtheorem{theorem}{Theorem}[section]
\newtheorem{lemma}[theorem]{Lemma}
\newtheorem{proposition}[theorem]{Proposition}
\newtheorem{corollary}[theorem]{Corollary}
\newtheorem{coro}[theorem]{Corollary}
\newtheorem{conj}[theorem]{Conjecture}
\theoremstyle{definition}
\newtheorem{definition}[theorem]{Definition}
\newtheorem{remark}[theorem]{Remark}
\newtheorem{rema}[theorem]{Remark}
\newtheorem{example}[theorem]{Example}
\newtheorem{question}[theorem]{Question}
\providecommand{\eg}{e.g.\ }
\providecommand{\ie}{i.e.\ }
\providecommand{\cf}{cf.\ }
\providecommand{\longthanks}[1]{\par\smallskip\noindent #1\par\smallskip}
\let\h\undefined
\numberwithin{equation}{section}
\newcommand \supp{\operatorname{supp}}
\newcommand \depth{\operatorname{depth}}
\newcommand \cd{\operatorname{codim}}
\newcommand \cochord{\operatorname{co-chord}}
\newcommand \reg{\operatorname{reg}}
\newcommand\bc{\operatorname{bc}}
\newcommand\link{\operatorname{link}}
\newcommand \Tor{\operatorname{Tor}}
\newcommand \ab{\operatorname{a}}
\newcommand \ca{\operatorname{c}}
\newcommand \ma{\operatorname{m}}
\newcommand \Vr{\operatorname{V}}
\newcommand \ba{\operatorname{b}}
\newcommand \h{\operatorname{ht}}
\newcommand \MG{\operatorname{MinGen}}
\newcommand \MM{\operatorname{min-match}}
\newcommand\n{\mathfrak{n}}
\newcommand\oddgirth{\operatorname{odd-girth}}
\newcommand \ka{\mathcal{K}}
\newcommand \cl{\operatorname{cl}}
\newcommand \C{\mathcal{C}}
\newcommand \G{\mathcal{G}}
\newcommand \NN{\mathbb{N}}
\newcommand \D{\mathcal{D}}
\newcommand\PP{\mathcal{P}}
\newcommand \DA{\mathcal{D}}
\newcommand \FA{\mathcal{F}}
\newcommand \A{\mathcal{A}}
\newcommand \B{\mathcal{B}}
%\newcommand \A{\mathcal{C}}
\newcommand \F{\mathcal{F}}
\newcommand \ha{\mathcal{H}}
\newcommand \St{\text{st}}
\newcommand \K{\mathbb{K}}
\newcommand \kk{\mathcal{K}}
\newcommand \s{\mathcal{S}}
\newcommand \pd{\operatorname{pd}}
\newcommand{\ass}{\operatorname{Ass}}
%% The title of the paper: amsart's syntax. 
\begin{document}
\section*{1269. Edge-neighbourhood completion across a fixed edge cannot increase regularity of a finite simple graph edge ideal over any field}
Banerjee, Arindam; Nevo, Eran. \emph{Regularity of Edge Ideals Via Suspension}. Algebraic Combinatorics 6 (2023), publisher version, DOI 10.5802/alco.317. \url{https://alco.centre-mersenne.org/articles/10.5802/alco.317/}. Exact official native/publisher PDF; publication posted2024-01-08. Native and PDF individually CC BY4.0, \url{https://creativecommons.org/licenses/by/4.0/}.
\paragraph{Selected problem (editorial).} Prove exactly Theorem3.1: for a finite simple graph G and an edge ab, adjoining every edge xy with x unequal y, ax and by edges, cannot increase the edge-ideal regularity over the original field. The full new complementary-clique-complex decomposition, minimal induced homology witness, both Mayer--Vietoris branches, induced-suspension transfer and antistar contradiction remain required core. The full other article body is retained as closed notation/reference context but receives no count.
\paragraph{Ordinary prerequisites (editorial).} Standard graded-module regularity, edge ideals and the cited Hochster homology formula, reduced-homology Mayer--Vietoris exact sequences and suspension identity are ordinary prerequisites. The new complex decomposition and transfer of the minimal witness are not assumed. Other named colon-ideal foundations in unselected later material remain cited prerequisites only.
\paragraph{Difficulty (editorial).} H2. The authors discover an exact decomposition of the complementary clique complex after neighbourhood completion and use a minimal induced homology witness, two Mayer-Vietoris cases, and an induced suspension to transfer nonvanishing homology back to the original graph. The homological obstruction is preserved through a nonlocal graph modification, rather than following from monotonicity of ordinary edge inclusion.
\paragraph{Literal source interfaces (editorial).} Original open-star versus closed-star notation, the literal final reg(Delta) wording at489, and the full matching graph/complex/regularity context are preserved. No notation repair, extra bridge or stronger assertion is inserted.
\paragraph{Formatting (editorial).} Complete exact native209--665, with the original macros49--86 and equation indexing27, is retained in source order including author comments, inactive commented material and all percent guards. Standalone article typography and front matter replace the publisher class; no source class, external dependency acquisition, private preamble or proof code is executed. All formulas and homology sequences are editable native TeX.
\clearpage
\section{Introduction}
Let $M$ be a finitely generated graded module over a polynomial ring $R = \K[x_1,\ldots,x_n]$, where
$\K$ is a field. The Castelnuovo--Mumford regularity (or simply, regularity) $\reg(M)$ of $M$
is defined as
$$
 \reg(M)=\max \{j-i \mid \Tor_i^R(M, \K)_j \neq 0\}.
$$
 $\text{  }$ $\text{  }$Regularity  is an important invariant in commutative algebra and algebraic geometry that measures in some sense the complexity of ideals, modules, and sheaves. A question that has been studied by many is how the regularity behaves with respect to taking powers of homogeneous ideals. It is known that in the long-run $\reg(I^k)$ is linear in $k$, that is, there exist integers $a(I), b(I), c(I)$ such that $\reg(I^k) = a(I)k + b(I)$ for all $k\geq c(I)$ (see \cite{CHT,vijay}). For various classes of ideals people have studied these integers and also have looked for various upper and lower bounds for $\reg(I^k)$.
  %For monomial ideals this study has a combinatorial angle to it
 %as they are equipped with nice combinatorial structures
 For monomial ideals these invariants, as well as bounds on them, reflect the underlying combinatorics (see e.g. \cite{BBH17,ha_adam2,Herzog'sBook,ms2005,MorVill,Nevo} for various works under this theme).
 For monomial ideals $I$ generated in same degree $d$, Kodiyalam~\cite{vijay} showed that $a(I)=d$.
 % Conca~\cite{C} showed that for any given integer $d > 1$ there exists an ideal $J$ generated by $d+5$ monomials of degree $d + 1$ in $4$ variables such that $\reg(J^k) = k(d + 1)$ for every $k <d$ and $\reg(J^d) \geq d(d + 1) + d-1$. In particular, $c(I)$ cannot be bounded above in terms of the number of variables only, not even for monomial ideals in general.

 One important class of monomial ideals is the class of edge ideals $I(G)$ of finite simple graphs, namely the ideals generated by squarefree monomials of degree two. For edge ideals, $c(I(G))\le 2$ for various cases: for example when the underlying graph is either cochordal or gap and cricket free or bipartite with $\reg(I(G))\leq 3$
 %We have $c(I(G))=1$ in for the cochordal graphs, $c(I(G)) \leq 2$ for gap and cricket free graphs and $c(I(G))=2$ for bipartite graphs with $\reg(I(G))\leq 3$
 (see \cite{banerjee1,banerjee,BBH17,froberg}). %For further results we refer to the references.
 As of $b(I(G))$, all known examples have
 \begin{equation}\label{eq:b(I)}
b(I(G))\leq \reg(I(G))-2
 \end{equation}
 and it is conjectured by Alilooee--Banerjee--Beyarslan--H\`a--Jayanthan--Selvaraja (see e.g. \cite{BBH17,JS18}) that this inequality holds for every graph.  In fact, the conjecture is
 %a slightly more general statement:
slightly stronger:
\begin{conj}\label{conj:byMany}\textup{(Alilooee--Banerjee--Beyarslan--H\`a--Jayanthan--Selvaraja)}
For every finite graph $G$, every $s\geq 1$ and every field $\K$,  $$\reg(I(G)^s)\leq 2s+ \reg(I(G))-2.$$
\end{conj}
Note that the values $\reg(I(G)^s)$ and even $b(I(G))$ may depend on the characteristic of the field, as was shown recently in~\cite[Exa. 1.2]{minh_vu}, and earlier for monomial ideals generated in higher degrees in~\cite{machia}.
For various classes of graphs (e.g. cochordal)
 %or bipartite with regularity $3$)
 Conjecture~\ref{conj:byMany} is known to be true.

  Our Theorem~\ref{thm:main} below proves this conjecture for all graphs for second powers and for all bipartite graphs for all powers. As a result it verifies inequality (\ref{eq:b(I)}) for all bipartite graphs.
 %In this paper we study the problem of finding upper bounds for powers of edge ideals. In particular we prove that for bipartite graphs $\reg(I(G))^k \leq 2k+\reg(I(G)) -2$, so indeed $b(I(G))\leq \reg(I(G))-2$.
 To demonstrate that this bound is sharp, we observe that all complete bipartite graphs $G$ with nonempty edge set satisfy $\reg(I(G)^s)=2s$ for all $s$, by the well known Theorem~\ref{thm:chordal} below.

 %and this achieves the value of our bound.

 Our main theorem is the following:
\begin{theorem}
\begin{enumerate}[label=(\roman*)]
\item \label{it:1} Let $G$ be a finite simple graph. Then
$$ \text{ }\reg(I(G)^2) \leq \reg (I(G))+2.$$
\item \label{it:2} Further, if $G$ is also bipartite, then for all $s\geq 1$ we have $$\text{ } \reg(I(G)^s) \leq 2s+\reg (I(G))-2.$$
\end{enumerate}
\label{thm:main}
\end{theorem}

Part~\ref{it:1} is proved topologically, via Hochster's formula and various uses of Mayer--Vietoris long exact sequence. Suspension plays a key role here. Part~\ref{it:2} for $s>2$ is proved algebraically, via various uses of short exact sequences for related ideals. Part~\ref{it:2} improves the main result of ~\cite{jayanthan}, which proves that if $G$ is bipartite, then for all $s\geq 2$  there holds $\reg(I(G)^s) \leq 2s+\text{cochord(G)}-1$, where cochord(G) is the cochordal number of $G$ (see \cite{jayanthan} for definition).

\textbf{Outline}: Preliminaries are given in Section~\ref{pre}, Theorem~\ref{thm:main} is proved in Section~\ref{sec:main_results}, and concluding remarks are given in Section~\ref{sec:further}.


\section{Preliminaries}\label{pre}
  In this section, we set up the basic definitions and notation needed for the main results. Let $M$ be a finitely generated graded $R=\K[x_1,\ldots,x_n]$-module.
Write the graded minimal free resolution of $M$ in the form:
\[
  0 \longrightarrow \bigoplus_{j \in \mathbb{Z}} R(-j)^{\beta_{p,j}(M)}
\overset{\psi_{p}}{\longrightarrow} \cdots \overset{\psi_1}{\longrightarrow}
\bigoplus_{j \in \mathbb{Z}} R(-j)^{\beta_{0,j}(M)}
\overset{\psi_0}{\longrightarrow} M\longrightarrow 0,
 \]
 where $p \leq n$, $R(-j)$ indicates the ring $R$ with the
shifted grading such that, for all $a \in \mathbb{Z}$, $R(-j)_{a}=R_{a-j}$.
The nonnegative integers $\beta_{(i,j)}(M)$ are called 
$i^{\text{th}}$-graded Betti numbers of $M$ in degree $j$.

The Castelnuovo-Mumford regularity (or regularity) of $M$ is defined to be
\[
 \reg(M)=\max\{j-i \mid \beta_{i,j}(M)\neq 0\}.\label{def_reg}
\]
  Let $I$ be a nonzero proper homogeneous ideal of $R$.
Then it follows from the definition
that $\reg(I) = \reg(R/I)+1.$

Let $I$ be any ideal of $R$ and $a\in R$ any element, the the \emph{colon ideal} $(I:a)$ is defined as the ideal $(I:a)\coloneqq (b \text{ } | b\in R, ab\in I)$.

    \emph{Polarization} is a process that creates a squarefree monomial out of a monomial, possibly in a larger polynomial ring. If $f=x_1^{e_1}\cdots x_n^{e_n}$ is a monomial in $\mathbb{K}[x_1, \ldots , x_n]$ then the polarization of $f$ is defined as $\tilde{f}=x_{11}\cdots x_{1e_1} x_{21} \cdots x_{2e_2}\cdots x_{n1} \cdots x_{n e_n}$ in the ring $\mathbb{K}[x_{11},\ldots x_{1e_1}, x_{21},\ldots x_{2e_2}, \ldots,  x_{n1}, \ldots, x_{ne_n}]$.
    For convenience
    we identify the variable $x_{i1}$ with $x_i$, so the new polynomial ring extends the old one.
    For a monomial ideal $I$ with minimal monomial generators $\{m_1,\ldots, m_k\}$, we define the \emph{polarization} of $I$ as $\widetilde{I}\coloneqq(\widetilde{m_1}, \cdots, \widetilde{m_k})$ in a suitable ring, see e.g
   \cite{Herzog'sBook} or \cite[Sec. 1.6]{kummini_thesis}. In the special case where degree of a variable $u=x_i$ is two in one or more generators then we call the unique new variable $x_{i2}$ a \emph{whisker variable} and denote it by $u'$ for short.
  %If $u$ is a variable whose degree is two  in the monomial that is being polarized then the new variable is denoted by $u'$.
  In this paper we repeatedly use an important property of 
polarization:
\begin{theorem}[e.g.~{\cite[Cor. 1.6.3(a)]{kummini_thesis}}] Let $I$ be a monomial
ideal in $R$.
%$\K[x_1, \ldots, x_n].$
Then
$$\reg(I)=\reg(\widetilde{I}).$$
\label{thm:2.1}
\end{theorem}

One of the main techniques that is used in this paper is that of short exact sequences. 
In particular we shall use the following well known result~\cite[Lemma. 2.11]{banerjee}:

\begin{theorem}
\begin{enumerate}[label=(\roman*)]
\item Let $I$ be a homogeneous ideal in a polynomial ring $R$ and $m$ be an element of degree $d$ in $R$. Then the following is a short exact sequence: $$0 \longrightarrow \frac{R}{(I:m)} \overset{\cdot m} \longrightarrow \frac{R}{I} \longrightarrow \frac{R}{I+(m)} \longrightarrow 0.$$
Hence $$\reg(I) \leq \max \{\reg (I:m) +d, \reg (I+(m))\}.$$
\item In case $I$ is squarefree and $x$ a variable, then also $\reg(I,x) \leq \reg I$.
\end{enumerate}
\label{thm:2.2}
\end{theorem}


 Let $G$ be a (finite simple) graph with vertex set $V(G)=\{x_1, \ldots ,x_n\}$ and the edge set $E(G)$. The  \textit{edge ideal} $I(G)$ of $G$ is defined as the ideal in $R$: $$I(G)=(x_i x_j|x_ix_j \in E(G)).$$
 For example, edge ideal of a $5$-cycle is $(x_1x_2,x_2x_3,x_3x_4,x_4x_5,x_5x_1)$.
 
 The next couple of theorems allow for induction when increasing the power of an edge ideal.
 
\begin{theorem}[{\cite[Thm. 5.2]{banerjee}}]
\label{thm:2.3}
For any graph $G$ and any $s \geq 1,$ let the set of minimal monomial generators of $I(G)^s$ be $\{m_1, \dots, m_k\}.$ Then
$$ \reg (I(G)^{s+1}) \leq \max \{\reg (I(G)^s) , \reg ((I(G)^{s+1}): m_l)+2s , 1 \leq l \leq k\}.$$
\end{theorem}

 By identifying the variables with the vertices of $G$, interpreting edges as squarefree quadratic monomials, defining neighborhood for any vertex $c$,  $N(c)\coloneqq\{z\in V(G): cz \in E(G)\}$ and using  \cite[Thm. 5.2]{banerjee}  we get the following corollary:

\begin{coro}
\begin{enumerate}[label=(\roman*)]
\item The ideal $(I(G)^{s+1}: m_l)$ is a quadratic monomial ideal.
\item For the special case where $s=1$ and $m=ab$ is an edge we have $$(I(G)^2:ab)=I(G)+(xy| x\in N(a),y\in N(b)).$$
\end{enumerate}
\label{cor:I2:e}
\end{coro}

 For bipartite graphs we further have:
\begin{theorem}[{\cite[Lem. 3.2, Prop. 3.4]{banerjee1}}]
Let $G$ be a bipartite graph and $s \geq 1$ an integer. 
Then $(I(G)^{s+1}: e_1 \cdots e_s)$ is a quadratic squarefree monomial ideal where for all $i$ we have $e_i \in E(G)$. Moreover, the graph $G'$ associated to $(I(G)^{s+1} : e_1\cdots e_s)$ is bipartite on the same vertex set and same bipartition as G.

 For any bipartite graph $G$ we have $$(I(G)^{s+1}:e_1\cdots e_s)=((I(G)^2:e_i)^s : \Pi_{j \neq i} e_j)$$ where for all $l \in \{1, \ldots,s\}$ we have $e_l \in E(G)$.
 \label{thm:2.5}
\end{theorem}

\begin{rema} From Theorem~\ref{thm:2.5} one can show a slightly more general equality  involving colons, under the same assumptions: for all $k$ and any $s \geq k+1$,
\begin{equation*}
(I(G)^{s+1}:e_1\cdots e_s)=((I(G)^{k+1}:e_1 \cdots e_k)^{s-k+1} : e_{k+1} \cdots e_s).
\end{equation*}
This follows from the following sequence of arguments:
\begin{enumerate}[align=parleft,left=0pt,label=(\alph*)]
\item From Theorem 2.5 for $i=1$ it follows that:
\begin{equation*}
  (I(G)^{s+1}:e_1\cdots e_s)=((I(G)^2:e_1)^s : e_2 \cdots e_s ).
\end{equation*}

 \item The first part of the Theorem~\ref{thm:2.5} tells us that $(I(G)^2:e_1)$ is a bipartite edge ideal. 
  \item By Corollary~\ref{cor:I2:e} above we get that each $e_i$ is an edge of the bipartite graph associated to $(I(G)^2:e_1)$.
  \item Applying the Theorem~\ref{thm:2.5} again on the edge ideal $(I(G)^2:e_1)$ we get:
 \begin{equation*}
 ((I(G)^2:e_1)^s : e_2 \cdots e_s )=(((I(G)^2:e_1)^2:e_2)^{s-1} : e_3 \cdots e_s).
 \end{equation*}
\item By a similar application of the last result with $s=2$ we have
  \begin{equation*}
  (I(G)^3:e_1e_2)=((I(G)^2:e_1)^2:e_2).
  \end{equation*}

\item Combining these we get
\begin{equation*}
 (I(G)^{s+1}:e_1\cdots e_s)=((I(G)^3:e_1e_2)^{s-1} : e_3 \cdots e_s).
 \end{equation*}

\item Inductively, assuming that for a fixed $k$ one has for any $s \geq k+1$
\begin{equation*}
(I(G)^{s+1}:e_1\cdots e_s)=((I(G)^{k+1}:e_1e_2 \cdots e_k)^{s-k+1} : e_{k+1} \cdots e_s).
\end{equation*}
\item Again applying Theorem~\ref{thm:2.5} on the bipartite edge ideal $(I(G)^{k+1}:e_1 \cdots e_k)$ we get
 \begin{equation*}
 ((I(G)^{k+1} :e_1 \cdots e_k)^{s-k+1} :e_{k+1} \cdots e_s)=( ((I(G)^{k+1}: e_1 \cdots e_k )^2:e_{k+1})^{s-k} : e_{k+2} \cdots e_s).
\end{equation*}
\item But by induction $$(I(G)^{k+2}:e_1 \cdots e_{k+1})=((I(G)^{k+1}: e_1 \cdots e_k )^2:e_{k+1}).$$
  So we get $$((I(G)^{k+1} :e_1 \cdots e_k)^{s-k+1} :e_{k+1} \cdots e_s)= ((I(G)^{k+2}:e_1 \cdots e_{k+1})^{s-k}:e_{k+2} \cdots e_s),$$ as claimed.
\end{enumerate}
\end{rema}

Now we recall some basic definitions about graphs and simplicial complexes that will be useful.

Let $G$ be a graph with vertex set $V(G)$
and edge set $E(G)$.  
A subgraph $H \subseteq G$  is called \textit{induced} if $\{u,v\}$ is an edge of $H$ if and only if $u$ and $v$ are vertices of $H$ and $\{u,v\}$ is an edge of $G$. 
A \textit{clique} in a graph is an induced subgraph that is a  \emph{complete graph}.

For $u\in V(G)$, let $N_G(u) = \{v \in V (G)\mid \{u, v\} \in E(G)\}$ and $N_G[u]=N_G(u)\cup \{u\}$.
For $U \subseteq V(G)$, denote by $G \setminus U$
the induced subgraph of $G$ on the vertex set $V(G) \setminus U$.

Let $G$ be a graph. We denote the graph consisting of two disjoint edges
 by $2K_2$.
%  in $G$ if $G$ does not have an edge with one
%endpoint in $f_1$ and the other in $f_2$.
A graph without induced copy of $2K_2$ is called $2K_2$-free or \textit{gap-free} graph. The \textit{complement} of a graph $G$, denoted by $G^c$, is the graph on the same
vertex set in which $\{u,v\}$ is an edge of $G^c$ if and only if it is not an edge of $G$.
Then  $G$ is gap-free if and only if $G^c$ contains no induced $4$-cycle.
%Thus, $G$ is gap-free if and only if it does not contain two vertex-disjoint edges as an induced subgraph.

%\par A \textit{matching} in a graph $G$ is a subgraph consisting of pairwise disjoint edges.
%If the
%subgraph is an induced subgraph, the matching is an \textit{induced matching}.
%T%he largest size of an induced matching in $G$ is called its induced
%matching number and denoted by $\nu(G)$.
%Let $H$ be a subgraph of $G$. The largest size of induced matching of $H$ as well as
%$G$ denoted by $\nu_{GH}$.  A \textit{dominating induced matching} of $G$ is an induced matching which also
%forms a maximal matching of $G$
A graph $G$ is chordal (also called triangulated) if every induced
cycle in $G$ has length $3$, and is co-chordal if the complement graph $G^c$ is chordal.
 The following important theorem(s) characterizes the edge ideals with regularity 2, and the regularity of their powers.

 

\begin{theorem}
\begin{enumerate}
\item \textup{(}Fr\"oberg~\cite[Thm.1]{froberg}\textup{)} For any graph $G$, we have $G^c$ is chordal if and only if $\reg(I(G))=2$. 

\item \textup{(}Herzog--Hibi--Zheng~\cite[Thm.1.2]{HHZ}\textup{)} Further, in this case $\reg(I(G)^s)=2s$ for any $s$.
\end{enumerate}
\label{thm:chordal}
\end{theorem}

  A \emph{simplicial complex} $\Delta$ on a vertex set $\{1,\ldots, n\}$ is a collection of subsets of $\{1,\ldots ,n\}$ such that if $\tau \in \Delta, \sigma \subseteq \tau$ then we have $\sigma \in \Delta$ and all singletons belong to that collection; if $\tau$ is such a subset belonging to $\Delta$ then $\tau$ is called a \emph{face} of $\Delta$. 
 The \emph{induced subcomplex} $\Delta [A]$ of $\Delta$ on vertex set $A\subset \{1,\ldots,n\}$ is the collection of faces $\tau'$ of $\Delta$ such that $\tau' \subset A$. 
 Clearly the induced subcomplex is a simplicial complex itself. We denote by $\Vr(\Delta)$ the set of vertices of $\Delta$.

 For sets  $A$ and $B$ we sometimes write $A\cup_C B$ for the set $A\cup B$ as a shorthand for denoting their intersection by $C=A\cap B$. 
 Likewise for simplicial complexes.

 The \emph{link} of a vertex $v$ in $ \Delta$ is $$\text{link}_v (\Delta)= \{ \tau |\tau \cup \{v\} \in \Delta, \{v\} \cap \tau =\emptyset\}.$$

 The (open) \emph{star} of a vertex $v$ in a simplicial complex $\Delta$ is the set of all faces that contain $v$, namely $\St_v(\Delta)=\{\tau| \tau \in \Delta, v \in \tau\}$. The \emph{closed star} $\bar{\St}_v (\Delta$ ) of $v$ is defined by the smallest subcomplex that contains $\St_v(\Delta)$. The \emph{antistar} of vertex $v$ is defined as the subcomplex $\text{ast}_v (\Delta) = \{\tau \in \Delta| \tau \cap \{v\} = \emptyset\}$.


 The \emph{join} of two simplicial complexes $\Delta_1$ and $\Delta_2$ is defined by $\Delta_1 * \Delta_2=\{\sigma \cup \tau| \sigma \in \Delta_1, \tau \in \Delta_2\}$. 
 The \emph{suspension} of a simplicial complex $\Delta$, w.r.t. two points $a$ and $b$ not vertices of $\Delta$, is the join defined by $\Sigma_{a,b} \Delta=\Delta * \{\{a\},\{b\},\emptyset \}$; its geometric realization is homeomorphic to the topological suspension of the space $\Delta$.

  For $H$ a graph
  %and $\Delta=\cl(G^c)$, where
  let $\cl(H)$ denote its clique complex, i.e. the simplicial complex whose faces are the cliques of $H$.

The following formulation of regularity follows from the so called Hochster's formula
%and Stanley-Reisner ideal formulation
(see {\cite{ms2005}} for further details):

\begin{theorem}[Hochster's formula]
For every graph $G$ whose edge set is nonempty and $\Delta=\cl(G^c)$ we have:
$$\reg(G)\coloneqq\reg(I(G))= \max\{l+2:\exists W \subseteq \Vr(\Delta), \widetilde{H}_l(\Delta[W];k)\neq 0\}.$$
\end{theorem}
(Sometimes Hochster's formula is phrased with the extra condition on $W$ that $|W|=l+2+i$ for some $i\ge 0$. However, this condition is superfluous: if $|W|\le l+1$ then either $\dim(\Delta[W])<l$ or $\Delta[W]$ is the simplex on $l+1$ vertices. In both cases  $\widetilde{H}_l(\Delta[W])=0$.)
%\iffalse
\section{Main Results}\label{sec:main_results}
 We first prove that $\reg(I(G)^2:e) \leq \reg(I(G))$ for every edge $e$ of a graph $G$. This will lead us to our main result via a series of short exact sequence arguments. For that we first prove the following:

\begin{theorem}\label{thm:reg(D')}
Let $G$ be a graph, $ab \in E(G)$ and $G'=G \cup \{xy: x\neq y,  ax,by \in E(G)\}$. Then
 $$\reg (I(G') ) \leq \reg(I(G)).$$
\end{theorem}

\begin{proof}
Denote $\Delta=\cl(G^c)$ and $\Delta'=\cl(G'^c)$.
Let $A=\Vr(\bar{\St}_a (\Delta)),B=\Vr(\bar{\St}_b(\Delta)), C=A\cap B$ and $D=\Vr(\Delta)-(A\cup B)$. With this notation we observe the following crucial relation between $\Delta$ and $\Delta'$:
\begin{equation}\label{eq:Delta'}
 \Delta'=(\Delta[A]\cup_{\Delta[C]} \Delta[B])\cup_{\Delta[C]} (\Delta[C] \cup_{d\in D} (\{d\}*\Delta[C \cap \Vr(\St_d (\Delta))])).
\end{equation}

 We obtain this equality simply from the definition of $G'$, where every neighbour of $a$ is connected to every neighbour of $b$.
 The above decomposition of $\Delta'$ shows us how to prove the theorem using Hochster's formula and Mayer--Vietoris sequences, detailed next.


 Let $\reg(I(G')) = l+2$ and $W$ be a subset of the vertices of minimal size such that $\widetilde{H}_l (\Delta'[W]) \neq 0$.
 Decompose $W=W_1 \cup_{W_C} W_2 \subseteq V$ where $W_1=W\cap(A\cup_C B)$ (recall $C=A\cap B$), $W_2=W\cap(C\cup D)$ and $W_C=W_1\cap W_2=C \cap W$.

If $W\cap D=\emptyset$ then Eq.(\ref{eq:Delta'}) implies $\Delta'[W]=\Delta[W]$ and thus $\reg(I(G))\ge l+2$ as claimed.
Thus, we may assume
$W\cap D\neq \emptyset$, so $|W_1|<|W|$ and thus, by minimality of $W$, $\widetilde{H}_l(\Delta'[W_1])=0$.

If $|W_2|<|W|$ then minimality of $|W|$ also imply $\widetilde{H}_l(\Delta'[W_2])=0$.
By~\eqref{eq:Delta'} the decomposition $\Delta'[W]=\Delta'[W_1] \cup_{\Delta'[W_C]} \Delta'[W_2]$ holds. Consider the following Mayer--Vietoris exact sequence: $$\widetilde{H}_l \Delta'[W_1] \bigoplus \widetilde{H}_l \Delta'[W_2] \longrightarrow \widetilde{H}_l \Delta'[W] \longrightarrow \widetilde{H}_{l-1} (\Delta'[W_C]).$$
As the two summands on the left term vanish, exactness implies
$\widetilde{H}_{l-1} (\Delta'[W_C])\neq 0$. 
Now by~\eqref{eq:Delta'} we have  $\Delta'[W_C]=\Delta[W_C]$, hence suspension gives
  $0 \neq \widetilde{H}_l(\Delta[W_C] * \{\{a\},\{b\},\emptyset\})= \widetilde{H}_l(\Delta[\{a,b\} \cup W_C])$.
 We conclude $\reg (\Delta) \geq l+2$ as claimed.

Thus we may assume additionally that
$W_2=W$. 
Denote $X=\Delta'[W]=\Delta'[W_2]$ and let $d\in W\cap D$.
By~\eqref{eq:Delta'} $\text{link}_d X$ is an induced subcomplex of $\Delta[C]$, and further, the suspension of $\text{link}_d X$ by the vertices $a$ and $b$ is an induced
subcomplex of $\Delta$.

Thus, if
$\widetilde{H}_{l-1}(\text{link}_d X)\neq 0$
then
$0\neq \widetilde{H}_l (\{\{a\},\{b\},\emptyset\}*\text{link}_d X)=\widetilde{H}_l(\Delta[\{a,b\}\cup V(\text{link}_d X)])$, and thus $\reg(I(G))\ge l+2$ as claimed. 
So we may assume $\widetilde{H}_{l-1}(\text{link}_d X)=0$.


Consider the Mayer--Vietoris long exact sequence
corresponding to the union
$X=\text{ast}_d X\cup_{\text{link}_d X} \bar{\text{St}}_d X$:
$$\widetilde{H}_l (\text{ast}_{d} X) \bigoplus \widetilde{H}_l (\overline{\text{St}}_{d} X) \longrightarrow \widetilde{H}_l (X) \longrightarrow \widetilde{H}_{l-1} (\text{link}_d X).$$
By assumption, the right term is zero while the middle term is nonzero, and closed stars have vanishing homology, hence
$\widetilde{H}_l (\text{ast}_{d} X)=
\widetilde{H}_l (\Delta'[W\setminus\{d\}])\neq 0$.
This contradicts the minimality of $W$, so in fact this last case cannot occur.
\end{proof}

\iffalse
\begin{remark}
The  essence of the above proof is that at least for one $X\in\{W \cap A, W \cap B, \{a,b\}\cup W_C, \Vr(\textnormal{link}_d \Delta) \cap W_2 \cap C, d \in D \}$homology groups $\widetilde{H}_l \Delta [X]$ is non zero.
\end{remark}
\fi

% \begin{theorem}
% % Let $I=J_0 \subseteq J_1 \cdots \subseteq J_n=J$ be a chain of edge ideals and
% % let $e_1\ldots e_n$ be an $n$-fold products of edges from $I$, for some $n \geq 1$.
% % Then $(J_0\cdots J_n : e_1\cdots e_n)$ is a quadratic monomial ideal generated by
% % the generators of $J$ and all quadratic monomial ideals $uv$ such that there exists a
% % sequence of variables $x_1, \ldots, x_{2k}$, for some $k \geq 1$, where the following holds:
% % \begin{enumerate}
% %  \item For all $0 \leq i \leq k-1$, $x_{2i+1} x_{2i+2}=e_j$, for some $1 \leq j \leq n$.
% %  \item For all $i$, $| \{l \geq 0 \mid x_{2l+1}x_{2l+2}=e_i\}| \leq |\{j \mid e_j=e_i\}|$.
% %  \item The quadratics $ux_1, x_2x_3 \ldots x_{2k} v$ are all generators of $J$ and for all $i$ at least $i$ of them belong to $J_i$.
% % \end{enumerate}
%   \end{theorem}
%
%  \begin{proof}
%  It is straght forward if $n=2$ (May be we can avoid this step?? If not we can copy paste what I wrote i your notebook).\\
%
%  Clearly each minimum generator has degree two. Then we proceed by induction. Suppose $uv$ is a minimum generator. If $uv$ is in $J$ then nothing to proof. Otherwise we have $f_0, \ldots, f_n$ with $f_i \in J_i$ and $f_0\ldots f_n$ divides $uv e_1 \ldots e_n$. So there exists a $i$ and $j$, say $i=1$ without loss of generality, such that $e_1=x_1 x_2$ and $ux_1$ is a generator of $J_j$. Hence $x_2v$ is a generator of $J_0 \ldots J_{j-1} J_{j+1} \ldots J_n : e_2\ldots e_n$. Hence by induction we are done.
%  \end{proof}

%\begin{remark}
 %If $G$ is triangle free then we notice that $\reg(I(\Delta'))= \reg(I(G)^2:ab)$
\iffalse
By corollary 2.3 and the fact that $\reg(\Delta)=\reg(I(G))$. Note by Theorem 2.3, Corollary 2.4 and Theorem 3.1 showed that:

\begin{corollary}
 For all graph $G$ we have $\reg(I(G)+(uv|u\neq v, u \in N(a),v\in N(b)) \leq \reg(I(G)).$
 \end{corollary}
\fi
 %We use this corollary to prove theorem 1.1 now.
We are now ready to prove Theorem~\ref{thm:main}.

%\end{remark}




\begin{proof}[Proof of Theorem~\ref{thm:main}]
\ref{it:1} By Theorem~\ref{thm:2.3} we just need to prove
\begin{equation}\label{eq:I(G)^2:ab}
\reg(I(G)^2:ab)\leq \reg(I(G))\ \ {\rm for\  all\ \ } ab \in E(G).
\end{equation}

Let $J=I(G) +(uv|u\neq v, u \in N(a),v\in N(b))$. By Corollary~\ref{cor:I2:e}
%2.4
we have:$$(I(G)^2:ab)=J+(u^2| u \in N(a) \cap N(b)).$$
By Theorem~\ref{thm:2.1}, $\reg(I(G)^2:ab)  = \reg(\widetilde{(I(G)^2:ab)})$, the polarization. Here $L\coloneqq\widetilde{(I(G)^2:ab)}=
%(I(G)^2:ab)=
J+(uu'| u \in N(a) \cap N(b))$
for new whisker variables $u'$
 in a larger polynomial ring (defined in Section~\ref{pre}).
So, it is enough to prove that $\reg(L) \leq \reg(I(G))$.


 If $N(a) \cap N(b)=\emptyset$ then the assertion follows by Theorem~\ref{thm:reg(D')}. Now let $N(a) \cap N(b)=\{u_1, \ldots, u_k\}$. Consider the following short exact sequences:
$$0 \longrightarrow \frac{R}{(L:u_1)} (-1) \rightarrow \frac{R}{L} \rightarrow \frac{R}{(L,u_1)} \rightarrow 0$$
$$0 \longrightarrow \frac{R}{((L, u_1):u_2)} (-1) \rightarrow \frac{R}{(L, u_1)} \rightarrow \frac{R}{(L,u_1,u_2)} \rightarrow 0$$
$$\vdots$$
$$0 \longrightarrow \frac{R}{((L,u_1,\ldots ,u_{k-1}):u_k)} (-1) \rightarrow \frac{R}{(L, u_1, \ldots , u_{k-1})} \rightarrow \frac{R}{(L,u_1,\ldots ,u_k)} \rightarrow 0.$$
  Now observe that $(L,u_1,\ldots ,u_k)=J+(\text{variables})$ and for every $i$,  $((L,u_1,\ldots u_{i-1}):u_i)=(L:u_i)+ (\text{variables})$. By repeated use of Theorem 2.2 (both parts) we have that $\reg (L) \leq \max \{ \reg (J), \reg(L:u)+1, u\in N(a) \cap N(b)\}$. Now by Theorem~\ref{thm:reg(D')} $\reg(J) \leq \reg(I(G))$. It is enough to show that $ \reg(L:u)+1 \leq \reg(I(G))$ for all $u \in N(a) \cap N(b)$.

 To prove this we use three facts:

 A. Going mod any set of variables  makes the regularity stay same or go down (follows from Theorem~\ref{thm:2.2} (ii)).

 B. Adjoining a disjoint set of variables keeps the regularity same  (follows directly from Hochster's formula (Theorem~\ref{thm:chordal})).

 C. Adding a disjoint edge makes the regularity go up by one. (This follows from Hochster's formula (Theorem~\ref{thm:chordal}) applied for the suspension of the original clique complex; e.g. by K\"unneth formula (see e.g. \cite[9.12]{ABJ}) or~\cite[Lemma.8]{russ}).

Now, $(L:u)=L+(u')+(s: u\neq s\in N_G(u)\cup N_G(a)\cup N_G(b))=I(G'[U])+(s: s\in N_{G'}(u))$ where $G'$ is the polarization graph with $L=I(G')$ and $U$ is the set of non-neighbours of $u$ in $G'$.
By Fact B, $\reg((L:u))=\reg(I(G'[U]))$.

Note that  $a,b\notin
U$, and that $G'[U]=G[U]$, as $u\in N_G(a)\cap N_G(b)$ and no edge involving a whisker vertex has its other vertex in $U$.
By Fact C,
$\reg(I(G[U])+(ab))=1+\reg(I(G[U]))$.

By Fact A, modding $I(G)$ by the variables in $(N(a)\cup N(b))\setminus \{a,b\}$, we obtain the inequality $\reg(I(G[U])+(ab))\leq \reg(I(G))$.
    
Combined, we have: $$\reg(L:u)+1=\reg(I(G[U]))+1\leq \reg(I(G)), $$
as desired. 
This proves~\eqref{eq:I(G)^2:ab} and hence proves part~\ref{it:1} of the theorem.

\ref{it:2} Here the underlying graph is bipartite. In part~\ref{it:1} we have already proved that $\reg(I(G)^2:e)\leq \reg(I(G))$ for any graph $G$ with $e \in E(G)$. We may assume $s \geq 3$ by part~\ref{it:1}.


By the second part of Theorem~\ref{thm:2.5} we get
$$ (I(G)^s:e_1\ldots e_{s-1}) =  ((I(G)^2:e_1)^{s-1}: e_2 \cdots e_{s-1})=  ((I(G)^2:e_1)^2 : e_2)^{s-2}: e_3 \cdots e_{s-1})$$
$$ = \ldots  =(\ldots((I(G)^2:e_1)^2 : e_2)^2 \ldots : e_{s-2})^2 : e_{s-1}).$$

Hence using the first part of Theorem~\ref{thm:2.5} and~\eqref{eq:I(G)^2:ab} that $\reg(I(G)^2:e)\leq \reg(I(G))$ for all graph $G$ and for all $e \in E(G)$ we get that:
$$ \reg(I(G)^s:e_1\ldots e_{s-1}) =
%\reg  (I(G)^2:e_1)^{s-1}: e_2 \cdots e_{s-1}=  \reg((I(G)^2:e_1)^2 : e_2)^{s-2}: e_3 \cdots e_{s-1})
\ldots =
\reg (\ldots((I(G)^2:e_1)^2 : e_2)^2 \ldots : e_{s-2})^2 : e_{s-1})$$
$$\leq \reg  (\ldots((I(G)^2:e_1)^2 : e_2)^2 \ldots : e_{s-1})$$
$$\leq \reg  (\ldots((I(G)^2:e_1)^2 : e_2)^2 \ldots : e_{s-2})\leq \ldots \leq \reg (I(G)^2:e_1)\leq \reg((I(G)).$$
This gives us the result by Theorem~\ref{thm:2.3}.
\end{proof}

\section{Further Research}\label{sec:further}
 In this section we discuss some questions for further research.
 
In the introduction we mentioned the constants $a(I), b(I), c(I)$ related to
asymptotic stability of the regularity of powers of a homogenous ideal $I$. So far we focused on $b(I)$ for edge ideals, now we turn to discuss $c(I)$.

 Due to the asymptotic stability we have that for a homogeneous ideal $I$ generated in degree $d$ there exists a minimal integer $c(I)$ such that $\reg(I^{s+1})-\reg(I^s)=d$ for all $s\geq c(I)$. 
 We have proved that for all bipartite graphs $G$ we have $\reg(I(G)^s)- \reg(I(G)) \leq 2s-2$. 
 However the behaviour of the sequence $\{\reg(I^s)\}$ can be irregular for smaller $s$ values even for edge ideals. In fact there are examples of bipartite graphs where $\reg(I(G)^2)=\reg(I(G)) +1$ (for example one can check that this is the case for the bipartite edge ideal of the 8-cycle  $(x_1y_1,x_2y_2,x_3y_3,x_4y_4,x_1y_2,x_2y_4,x_3y_1,x_4y_3)$).

 Can $c(I)$ be bounded by some simple invariants of $I$, for homogenous ideals?
  Conca~\cite{Conca} showed that for any given integer $d > 1$ there exists an ideal $J$ generated by $d+5$ monomials of degree $d + 1$ in $4$ variables such that $\reg(J^k) = k(d + 1)$ for every $k <d$ and $\reg(J^d) \geq d(d + 1) + d-1$. In particular, $c(I)$ cannot be bounded above in terms of the number of variables only, not even for monomial ideals in general. Further, a result of Raicu~\cite{Raicu} gives binomial ideals $I_n$ on $n^2$ variables, generated in degree $2$, with $c(I_n)=n-1$. All these fit into the framework of the following question that served as broad aim for various researchers
since the works of Cutkosky, Herzog and Trung, and Kodiyalam (\cite{BBH17,Conca,CH03,CHT, huneke,ha_adam2,HHZ, JS18,vijay,Nevo, nevo_peeva}):

\begin{question}
For homogeneous ideals $I$ on $n$ variables, generated in degree $d$, is $c(I)$ bounded above by a function of $d$ and $n$?
\end{question}

 It has been conjectured by Banerjee and Mukundan~\cite{arindam_vivek} that for all bipartite graphs $G$, we have $c(I(G)) \leq 2$. 
 It is known  for cochordal, gap free plus cricket/diamond/$4$-cycle free graphs~\cite{banerjee,HHZ,Er1,Er2}. Apart from edge ideals, it was shown by Conca and Herzog~\cite{CH03} that polymatroidal ideals have linear resolutions and powers of polymatroidal ideals are polymatroidal ideals. So for the class of all polymatroidal ideals $c(I)=1$.

Finally we conclude by a discussion on a related conjecture by~\cite{nevo_peeva}:

\begin{conj}[{\cite{nevo_peeva}}]\label{conj:Nevo-Peeva}
If $\reg(I(G)) \leq 3$ and $G^c$ has no induced $4$-cycle then for all $s \geq 2$ we have $\reg (I(G)^s) =2s$.
\end{conj}
Very recently this conjecture was verified for $s=2$ and $s=3$ by Minh and Vu~\cite{minh_vu2}.
 Theorem~\ref{thm:2.3} was proved by Banerjee in his thesis to study this conjecture and related other problems, based on Theorem~\ref{thm:2.2}.
We now explain why this inductive approach via colon ideals that is used in this paper (and also in \cite{banerjee1,banerjee,BBH17,Nevo,nevo_peeva,jayanthan,JS18}) can \emph{not} be used directly to settle Conjecture~\ref{conj:Nevo-Peeva}.

  Every $2$-dimensional simplicial complex $\Delta$ can be subdivided so that the resulted complex is flag-no-square, see~\cite{Dr} or~\cite[Lem. 2.3]{PS}, i.e. $\Delta=\text{cl}(H)$ where $H$ is a graph with no \emph{induced} $4$-cycles. In particular,
we choose such $H$ so that $\Delta$ triangulates the dunce hat, a contractible $2$-dimensional complex.
Thus, all subcomplexes of $\Delta$ have vanishing homology in dimension $\ge 2$. Further, the link of every vertex $a\in \Delta$ is an induced subcomplex (as $\Delta$ is a clique complex) with nonzero first homology. For an edge $ab\in G\coloneqq H^c$, the construction of $\Delta'=\cl(G'^c)$ from  Theorem~\ref{thm:reg(D')} satisfies $\text{link}_a \Delta'=\text{link}_a \Delta$ is an induced subcomplex of $\Delta'$.

We conclude that $\reg(I(G))=3$ and by Corollary~\ref{cor:I2:e}(ii)
for every edge $ab\in G$ also $\reg((I(G)^2:ab)=3$.
Thus, if $\reg(I(G)^2)=4$ as Conjecture~\ref{conj:Nevo-Peeva} suggests, then Theorem~\ref{thm:2.2} cannot be directly applied to prove it.

On the other hand, if $\reg(I(G)^2)>4$ then this will be a counterexample. Unfortunately we could not verify the value of $\reg(I(G)^2)$ due to computational limitations. It will be great if this can be verified in the future.


\longthanks{We thank Aldo Conca for pointing us to~\cite{Raicu}, and the referees for very helpful remarks on the presentation.

Most of this work was done when first author was visiting the Institute of Mathematics of Hebrew University and he would like to thank faculty and staff of Hebrew University for their hospitality. Also, the first author was partially supported by DST INSPIRE (India) research grant (DST/INSPIRE/04/2017/000752) and he would like to acknowledge that.

The second author was partially supported by Israel Science Foundation grants ISF-1695/15 and ISF-2480/20, by grant 2528/16 of the ISF-NRF Singapore joint research program, and by ISF-BSF joint grant 2016288.}

\begin{thebibliography}{99}
\bibitem{banerjee1}Alilooee, Ali; Banerjee, Arindam Powers of edge ideals of regularity three bipartite graphs, J. Commut. Algebra, Volume 9 (2017) no. 4, pp. 441-454
\bibitem{banerjee}Banerjee, Arindam The regularity of powers of edge ideals, J. Algebraic Combin., Volume 41 (2015) no. 2, pp. 303-321
\bibitem{BBH17}Banerjee, Arindam; Beyarslan, Selvi Kara; Hà, Huy Tài Regularity of powers of edge ideals: from local properties to global bounds, Algebr. Comb., Volume 3 (2020) no. 4, pp. 839-854
\bibitem{arindam_vivek}Banerjee, Arindam; Mukundan, Vivek The powers of unmixed bipartite edge ideals, J. Algebra and Its Appl., Volume 26 (2019), pp. 57-70
\bibitem{ABJ}Björner, A. Topological methods, Handbook of combinatorics, Elsevier, Amsterdam, Volume 1,2 (1995) no. 1, pp. 1819-1872
\bibitem{machia}Bolognini, D.; Machia, A.; Strazzanti, F.; Welker, V. Powers of Monomial Ideals With Characteristic-Dependent Betti Numbers, Research in the Mathematical Sciences, Volume 9 (2022) no. 26, pp. 630-645
\bibitem{Conca}Conca, Aldo Regularity jumps for powers of ideals, Commutative algebra (Lect. Notes Pure Appl. Math.), Volume 244, Chapman \& Hall/CRC, Boca Raton, FL, 2006, pp. 21-32
\bibitem{CH03}Conca, Aldo; Herzog, Jürgen Castelnuovo-Mumford regularity of products of ideals, Collect. Math., Volume 54 (2003) no. 2, pp. 137-152
\bibitem{CHT}Cutkosky, S. Dale; Herzog, Jürgen; Trung, Ngô Viêt Asymptotic behaviour of the Castelnuovo-Mumford regularity, Compositio Math., Volume 118 (1999) no. 3, pp. 243-261
\bibitem{huneke}Dao, Hailong; Huneke, Craig; Schweig, Jay Bounds on the regularity and projective dimension of ideals associated to graphs, J. Algebraic Combin., Volume 38 (2013) no. 1, pp. 37-55
\bibitem{Dr}Dranishnikov, A.N. Boundaries of Coxeter Groups and Simplicial Complexes with Given Links, J. Pure and Appl. Alg., Volume 137 (1999) no. 1, pp. 139-151
\bibitem{Er2}Erey, N. Powers of edge ideals with linear resolutions, Comm. in Alg., Volume 46 (2018) no. 9, pp. 4007-4020
\bibitem{Er1}Erey, N. Powers of Ideals Associated To (C 4 ,2K 2 )-Free Graphs, J. Pure Appl. Algebra, Volume 223 (2019) no. 7, pp. 3071-3080
\bibitem{froberg}Fröberg, Ralf On Stanley-Reisner rings, Topics in algebra, Part 2 (Warsaw, 1988) (Banach Center Publ.), Volume 26, PWN, Warsaw, 1990, pp. 57-70
\bibitem{ha_adam2}Hà, Huy Tài; Van Tuyl, Adam Resolution of Squarefree Monomial Ideals Via Facet Ideals: A Survey, Contemporary Mathematics, Volume 448 (2007) no. 2, pp. 91-117
\bibitem{Herzog'sBook}Herzog, Jürgen; Hibi, Takayuki Monomial ideals, Graduate Texts in Mathematics, 260, Springer-Verlag London, Ltd., London, 2011, xvi+305 pages
\bibitem{HHZ}Herzog, Jürgen; Hibi, Takayuki; Zheng, Xinxian Monomial ideals whose powers have a linear resolution, Math. Scand., Volume 95 (2004) no. 1, pp. 23-32
\bibitem{jayanthan}Jayanthan, A. V.; Narayanan, N.; Selvaraja, S. Regularity of powers of bipartite graphs, J. Algebraic Combin., Volume 47 (2018) no. 1, pp. 17-38
\bibitem{JS18}Jayanthan, A. V.; Selvaraja, S. Upper bounds for the regularity of powers of edge ideals of graphs, J. Algebra, Volume 574 (2021), pp. 184-205
\bibitem{vijay}Kodiyalam, Vijay Asymptotic behaviour of Castelnuovo-Mumford regularity, Proc. Amer. Math. Soc., Volume 128 (2000) no. 2, pp. 407-411
\bibitem{kummini_thesis}Kummini, Manoj Homological invariants of monomial and binomial ideals, Ph. D. Thesis, University of Kansas (2008)
\bibitem{ms2005}Miller, Ezra; Sturmfels, Bernd Combinatorial commutative algebra, Graduate Texts in Mathematics, 227, Springer-Verlag, New York, 2005, xiv+417 pages
\bibitem{minh_vu}Minh, C.; Vu, T. Integral closure of powers of edge ideals and their regularity., J. Algebra, Volume 609 (2022) no. 5, pp. 120-144
\bibitem{minh_vu2}Minh, C.; Vu, T. Characterization of graphs whose a small power of their edge ideals has a linear free resolution, Combinatorica (2023)
\bibitem{MorVill}Morey, Susan; Villarreal, Rafael H. Edge ideals: algebraic and combinatorial properties, Progress in commutative algebra 1, de Gruyter, Berlin, 2012, pp. 85-126
\bibitem{Nevo}Nevo, Eran Regularity of edge ideals of C 4 -free graphs via the topology of the lcm-lattice, J. Combin. Theory Ser. A, Volume 118 (2011) no. 2, pp. 491-501
\bibitem{nevo_peeva}Nevo, Eran; Peeva, Irena C 4 -free edge ideals, J. Algebraic Combin., Volume 37 (2013) no. 2, pp. 243-248
\bibitem{PS}Przytycki, P.; Swiatkowski, J. Flag-No-Square Triangulations And Gromov Boundaries in Dimension 3, Groups Geom. Dyn., Volume 3 (2009) no. 1, pp. 453-468
\bibitem{Raicu}Raicu, Claudiu Regularity and cohomology of determinantal thickenings, Proc. Lond. Math. Soc. (3), Volume 116 (2018) no. 2, pp. 248-280
\bibitem{russ}Woodroofe, Russ Matchings, coverings, and Castelnuovo-Mumford regularity, J. Commut. Algebra, Volume 6 (2014) no. 2, pp. 287-304
\end{thebibliography}
\end{document}
